% GitHub Repo and Documentation: https://github.com/celiobjunior/resume-template
% Copyright © 2025 Celio B Junior. All rights reserved.
% 
% Licensed under the Apache License, Version 2.0 (the "License");
% you may not use this file except in compliance with the License.
% You may obtain a copy of the License at
%
%     http://www.apache.org/licenses/LICENSE-2.0
%
% This template follows best practices from README.md

% Início do documento LaTeX: define tipo e formato do documento
% a4paper = tamanho A4, 10pt = tamanho base da fonte
\documentclass[a4paper,10pt]{article}

% --- PACOTES: BASE / IDIOMA ---
\usepackage[utf8]{inputenc}
\usepackage[T1]{fontenc}
\usepackage[portuguese]{babel}

% --- PACOTES: LAYOUT / TIPOGRAFIA ---
\usepackage{geometry}
\usepackage{parskip}
\usepackage{microtype}
\usepackage{enumitem}
\usepackage{titlesec}
\usepackage{array}
\usepackage{tabularx}

% --- PACOTES: LINKS / BOOKMARKS ---
\usepackage{hyperref}
\usepackage{bookmark}
\usepackage{xurl}

% --- CONFIGURAÇÃO: MARGENS / ESPAÇAMENTO ---
\geometry{top=1.5cm, bottom=1.5cm, left=1.5cm, right=1.5cm} % Margens da página
\setcounter{secnumdepth}{0}                                 % Remove numeração de seção
\setlist[itemize]{
    leftmargin=0.75em, % Recuo da lista (mais perto da margem da página)
    itemsep=0.2em,     % Espaço entre itens
    topsep=0.25em,     % Espaço antes/depois da lista
    parsep=0em,        % Espaço entre parágrafos no mesmo item
    partopsep=0em      % Espaço extra ao iniciar a lista
}

% Remove números de página e cabeçalhos para visual limpo do currículo
\pagestyle{empty}

% --- CONFIGURAÇÃO: METADADOS / LINKS ---
\hypersetup{
    pdftitle={CV Thiago Macedo Mendes},   % Título do PDF
    pdfauthor={Thiago Macedo Mendes},     % Autor do PDF
    colorlinks=true,          % Usa cor nos links (sem caixa)
    linkcolor=black,          % Cor de links internos
    urlcolor=black,           % Cor de URLs
    citecolor=black,          % Cor de citações
    bookmarksdepth=2          % Barra lateral: seção e subseção
}

% --- CONFIGURAÇÃO: TÍTULOS ---
\titleformat{\section}
{\Large\bfseries}
{}
{0em}
{}
[\titlerule\vspace{0.5ex}]
\titleformat{\subsection}
{\normalsize\bfseries}
{}
{0em}
{}
\titlespacing*{\subsection}{0pt}{0.35em}{0.15em}

% Macro para entradas do currículo com bookmark no PDF
\newcounter{cventry}
\newcommand{\cventry}[4]{%
    \refstepcounter{cventry}%
    \phantomsection%
    \pdfbookmark[2]{#1}{cventry-\thecventry}%
    \noindent\begin{tabularx}{\textwidth}{@{}>{\raggedright\arraybackslash}X >{\raggedleft\arraybackslash}X@{}}
    \textbf{#1} & #2 \\
    \textit{#3} & \textit{#4}
    \end{tabularx}\par
}

% --- INÍCIO DO DOCUMENTO ---
\begin{document}

% --- CABEÇALHO ---
% Substitua por suas informações pessoais
\begin{center}
    {\LARGE \textbf{Thiago Macedo Mendes}} 
    \\ [0.1cm]
    Porto Velho, Rondônia, Brasil
    {\textbullet}
    +55 (69) 99314-6868
    {\textbullet}
    \href{mailto:thiagomm@proton.me}{thiagomm@proton.me} 
    {\textbullet}
    \href{https://linkedin.com/in/thiagomacedomendes}{linkedin.com/in/thiagomacedomendes} 
    {\textbullet}
    \href{https://github.com/o-thiago}{github.com/o-thiago}
\end{center}

% --- RESUMO ---
\section{Resumo}
Estudante de Ciência da Computação (UNIR) e Técnico em Informática (IFRO) com atuação em desenvolvimento de software, infraestrutura de sistemas e pesquisa acadêmica. Experiência no desenvolvimento de aplicações full-stack e sistemas com TypeScript, Python, Java, Rust, PHP e outras, além de frameworks modernos (React, Next.js). Habilidade na padronização de ambientes reproduzíveis (Docker, Nix) e integração de ferramentas modernas de desenvolvimento com IA (Antigravity CLI, Zed AI, Codex). Trajetória com pesquisas em mecatrônica e sistemas web, premiações em olimpíadas científicas e apresentações técnicas em eventos públicos.

% --- EXPERIÊNCIA ---
\section{Experiência}
    % Liste suas experiências profissionais ou acadêmicas, da mais recente para a mais antiga
    \cventry{IFRO - Instituto Federal de Rondônia}{Porto Velho, Rondônia, Brasil}{Assistente Administrativo}{Jan 2025 - Dez 2025}
        \begin{itemize}
            \item Desenvolvi aplicações web full-stack (Next.js, React, Tailwind CSS, oRPC) para fluxos de pesquisa internos do grupo GPMecatrônica (em parceria com a Belém Rio Segurança Ltda).
            \item Automatizei e padronizei ambientes de desenvolvimento com Docker e Nix, garantindo reprodutibilidade e acelerando o onboarding de pesquisadores.
            \item Colaborei com equipes multidisciplinares em iniciativas de robótica e na coautoria de artigos científicos voltados a soluções tecnológicas aplicadas.
            \item Prestei suporte técnico avançado, diagnósticos de hardware e garantia de disponibilidade da infraestrutura do laboratório.
        \end{itemize}

    \cventry{Centro Gestor e Operacional do Sistema de Proteção da Amazônia (Censipam)}{Porto Velho, Rondônia, Brasil}{Estagiário de TI}{Ago 2023 - Out 2024}
        \begin{itemize}
            \item Atuei na sustentação de infraestrutura crítica de TI e conectividade de rede no Censipam (Centro Gestor e Operacional do Sistema de Proteção da Amazônia, Ministério da Defesa).
            \item Configurei pontos de telefonia IP e infraestrutura de rede estruturada para assegurar a continuidade da comunicação operacional.
            \item Realizei manutenção preventiva de hardware, diagnósticos e implantação de sistemas operacionais em conformidade com normas institucionais de segurança e proteção de dados.
        \end{itemize}

    \cventry{IFRO - Instituto Federal de Rondônia}{Porto Velho, Rondônia, Brasil}{Pesquisador Científico}{Fev 2024 - Jul 2024}
        \begin{itemize}
            \item Desenvolvi funcionalidades front-end e back-end para a plataforma web "Radar da Inovação" utilizando TypeScript, React, Laravel, PHP e MySQL no grupo de pesquisa GoTec.
            \item Estruturei biblioteca de componentes reutilizáveis documentada no Storybook, acelerando ciclos de desenvolvimento e assegurando consistência de design.
            \item Implementei conteinerização com Docker para garantir paridade entre ambientes de desenvolvimento local e produção.
        \end{itemize}

    \cventry{IFRO - Instituto Federal de Rondônia}{Porto Velho, Rondônia, Brasil}{Pesquisador Voluntário}{Set 2023 - Fev 2024}
        \begin{itemize}
            \item Representei o grupo de pesquisa GoTec em eventos de difusão científica, realizando apresentações na Expo Favela sobre tecnologia e inovação.
            \item Participei da concepção da arquitetura inicial de sistemas e modelagem de dados do projeto "Radar da Inovação".
            \item Conduzi treinamentos técnicos em desenvolvimento web (TypeScript, React, Laravel) para novos pesquisadores e estudantes integrados.
        \end{itemize}

% --- EDUCAÇÃO ---
\section{Educação}
    % Formação mais recente primeiro
    \cventry{Universidade Federal de Rondônia}{Porto Velho, Rondônia, Brasil}{Bacharelado, Ciência da Computação}{Fev 2026 - Dez 2029}\vspace{0.35em}
    \cventry{IFRO - Instituto Federal de Rondônia}{Porto Velho, Rondônia, Brasil}{Técnico em Informática Integrado ao Ensino Médio}{Fev 2023 - Dez 2025}

% --- HABILIDADES ---
\section{Habilidades}
    \begin{itemize}
        \item \textbf{IA e Ferramentas de Desenvolvimento:} Engenharia de Prompt, Fluxos Agênticos, Model Context Protocol (MCP), Antigravity CLI, Zed AI, Codex, Desenvolvimento Assistido por IA, Docker, Nix, Git, Linux.
        \item \textbf{Linguagens e Frameworks:} TypeScript, Python, React, Next.js, Java, Rust, C, PHP, Laravel, Tailwind CSS, oRPC.
        \item \textbf{Bancos de Dados e Infraestrutura:} PostgreSQL, MySQL, MongoDB, Infraestrutura de TI, Redes de Computadores.
        \item \textbf{Idiomas:} Português (Nativo ou Bilíngue), Inglês (Nativo ou Bilíngue), Espanhol (Básico).
    \end{itemize}

% --- CERTIFICAÇÕES ---
\section{Certificações}
    \begin{itemize}
        \item Medalha de Ouro - Olimpíada do Tesouro Direto de Educação Financeira 2025
        \item Medalha de Prata - Olimpíada Brasileira de Robótica (Teórica) 2025
        \item Menção Honrosa - Olimpíada Internacional de Matemática sem Fronteiras 2025
        \item Honra ao Mérito - Olimpíada Brasileira de Robótica (Teórica) 2024
        \item Participação na Olimpíada Brasileira de Informática (P2 em 2024, 2025; P1 em 2023)
        \item Palestrante no 1° Encontro Poliglota do Calama - Fundamentos do Inglês (2024)
        \item Bootcamps Desenvolvedor(a) Python e Full Stack (IGTI)
        \item Curso "Oratória: a arte de se expressar bem" (IFRO 2024)
        \item Curso de Tratamento e Proteção de Dados Pessoais (Ministério da Defesa 2024)
        \item Expositor na Semana Rondoniense de Arte Patrimônio e Cultura 2023
        \item Participação na Olimpíada Brasileira de Astronomia e Astronáutica 2025
        \item Participação na Semana Nacional de Ciência e Tecnologia (2023, 2024)
        \item Participação no Festival de Cultura Identidade e Educação Ambiental 2024
        \item Oficina de Elaboração de Currículo e Preparação para Entrevistas 2025
    \end{itemize}

\end{document}